In einem Inertialsystem gilt das Trägheitsgesetz: Ein Körper, auf den keine resultierende Kraft wirkt, bleibt in Ruhe oder bewegt sich geradlinig gleichförmig.

Bei einer Vollbremsung beginnen sich aber lose Gegenstände im Auto (z.B. eine Tasche auf dem Rücksitz) und die Insassen relativ zum Auto nach vorne zu bewegen, obwohl kein anderer Körper sie nach vorne zieht oder stösst. Die Insassen werden in die Gurte gedrückt. Im Bezugssystem Auto werden also Körper ohne Kraft beschleunigt: Das Trägheitsgesetz gilt nicht, das Auto ist $\result{\text{kein Inertialsystem}}$.

Von der Strasse aus betrachtet ist der Vorgang einfach: Die Gegenstände bewegen sich wegen ihrer Trägheit mit der ursprünglichen Geschwindigkeit weiter, während das Auto abgebremst wird.
